%document class
\documentclass[10pt,oneside]{book}
%%%% Page Info + Commands %%%%%{

%packages
\usepackage{pgfplots}
\pgfplotsset{compat=1.18}
\usepackage{amsmath}
\usepackage{amsfonts}
\usepackage{charter}
\usepackage{amsthm,letltxmacro}
\usepackage{amssymb}
\usepackage{fancyhdr}
\usepackage{float}
\usepackage[most]{tcolorbox}
\usepackage{enumerate}
\usepackage{enumitem}
\usepackage[dvipsnames]{xcolor}
\usepackage{changepage}
\usepackage{mdframed}
\usepackage{graphicx}

% SMALL PAGE COMMAND
\newcommand{\smallpage}{  \setlength \oddsidemargin{.25in}
            \setlength \textwidth{5.5in}
            \setlength \topmargin{0in}
            \setlength \textheight{9in}}


% Commands for sets of numbers
\newcommand{\Z}{\mathbb{Z}}\newcommand{\R}{\mathbb{R}}\newcommand{\N}{\mathbb{N}}\newcommand{\Q}{\mathbb{Q}}\newcommand{\C}{\mathbb{C}}
\newcommand{\real}[1]{\text{Re}\left(#1\right)}
\newcommand{\im}[1]{\text{Im}\left(#1\right)}

% gordie spacing and boxing commands
\newcommand{\gspc}{\vspace{0.13cm}} % 0.5 space
\newcommand{\gline}[0]{\gspc\\} % 1.5 space
\newcommand{\gbline}{\gspc\gspc\\} % double space
\newcommand{\gbox}[1]{\fbox{\parbox{\linewidth}{#1}}} % multiline box command
\newcommand{\gmod}[1]{\,(\text{mod}\,#1)}


\newcommand{\gimage}[2]{
\begin{figure}[H]
    \centering
    \includegraphics[width=#2]{#1}
\end{figure}
}

\newcommand{\centerit}[1]{{\begin{center} #1 \end{center}}}
\newcommand{\ul}[1]{\underline{#1}}


% COLOR BOX
\newmdenv[backgroundcolor=gray!8,
  linewidth=0pt, innerleftmargin=0pt, innerrightmargin=0pt,
  skipabove=\baselineskip, skipbelow=\baselineskip
]{coloredadjust}

% PROOF
\LetLtxMacro\oldproof\proof
\let\endoldproof\endproof
\renewenvironment{proof}[1][\proofname]
  {\vspace{0.2cm}\begin{adjustwidth}{2em}{2em}
   \oldproof[#1]}
  {\endoldproof
\end{adjustwidth}}

% PROBLEM
\newenvironment{problem}[1]
  {\vspace{-0.1cm}\begin{coloredadjust}\begin{adjustwidth}{2em}{2em}
   \textit{Problem. }#1}
  {\endoldproof
\end{adjustwidth}\end{coloredadjust}\gspc}

%% DEFINITION
\newtcbtheorem[]{definition}{Definition}{enhanced,
  colback=blue!2, colframe=blue!50!black, coltitle=blue!70!black,
  fonttitle=\bfseries,
  boxrule=0pt, toprule=2pt, bottomrule=1pt, leftrule=1pt, rightrule=1pt, arc=1pt, outer arc=1pt,
  attach title to upper, title={#1},
}{dfn}

%% CRITERION
\newtcbtheorem[]{criterion}{Criterion}{enhanced,
  colback=magenta!3, colframe=magenta!60!black, coltitle=magenta!20!black,
  fonttitle=\bfseries,
  boxrule=4pt, toprule=2pt, bottomrule=1pt, leftrule=1pt, rightrule=1pt, arc=1pt, outer arc=1pt,
  attach title to upper, title={#1},
}{ctn}

%% COROLLARY
\newtcbtheorem[]{corollary}{Corollary}{enhanced,
  colback=orange!3, colframe=orange!80!black, coltitle=orange!50!black,
  fonttitle=\bfseries,
  boxrule=4pt, toprule=1pt, bottomrule=1pt, leftrule=1pt, rightrule=1pt, arc=5pt, outer arc=5pt,
  attach title to upper, title={#1},
}{ctn}

%% THEOREM
\newtcbtheorem[]{theorem}{Theorem}{enhanced,
  colback=green!2, colframe=green!40!black, coltitle=green!20!black,
  fonttitle=\bfseries,
  boxrule=4pt, toprule=3pt, bottomrule=1pt, leftrule=1pt, rightrule=1pt,arc=5pt, outer arc=5pt,
  attach title to upper, title={#1},
}{thm}

%%%% Page Setup %%%%%%%
\smallpage\sloppy
\pagestyle{fancy}
\lhead{\large{\textbf{Feb 9th{}-{}Notes}}}
%\rfoot{\thepage/\pageref{LastPage} }
\setlength{\headheight}{10pt}

%%%%%%%%%%%%% PHYSICS SPECIFIC COMMANDS

\newcommand{\vA}{\vec{\mathbf{A}}}
\newcommand{\vB}{\vec{\mathbf{B}}}
\newcommand{\vC}{\vec{\mathbf{C}}}
\newcommand{\vZ}{\vec{\mathbf{0}}}
\newcommand{\vv}{\vec{\mathbf{v}}}
\newcommand{\norm}{\vec{\mathbf{n}}}
\newcommand{\pvec}[1]{\left\langle#1\right\rangle}
\newcommand{\ar}{\vec{\mathbf{r}}}
\newcommand{\arhat}{\hat{\mathbf{r}}}
\newcommand{\prt}[2]{\frac{\partial#1}{\partial#2}}
\newcommand{\prtd}[2]{\frac{\partial^2#1}{\partial{#2}^2}}
\newcommand{\cpr}[1]{\left(#1\right)}

%%%%%%%%%%%%%%%%%%%%%%%%%%%%
%%%%%%%%%%%%%%%%%%%%%%%%%%%%%%
%%%%%%%%%%%%%%%%%%%%%%%%%%%%%
%%%%%%%%%%%%%%%%%%%%%%%%%%%%%
%%%%%%%%%%%%%%%%%%%%%%%%%%%%%%
%%%%%%%%% begin doc


\begin{document}
\subsection*{Gradients}
Consider our Del operator:
\begin{align*}
  \nabla = \hat x \prt{}{x} + \hat y \prt{}{y} + \hat z\prt{}{z}
\end{align*}
When we take the Del of a scalar function, we call that the \textbf{gradient}. 
\begin{align*}
  \nabla T = \prt{T}{x}\hat x + \prt{T}{y}\hat y+\prt{T}{z}\hat z
\end{align*}
This gradient points in the direction of maximum change. Magnitude is the rate of change in that direction. \gline{}
However, something interesting occurs when we look at the separation vector, $\mathbf r$. Let $\mathbf r$ be the separation vector with fixed point $\pvec{x', y', z'}$ to point $\pvec{x, y, z}$. We want to show that:
\begin{align*}
  \nabla(1/\ar) = -\arhat/r^3
\end{align*}
The proof is left as an exercise to the reader. 
\subsection*{Divergence}
Consider the formula for gradience:
\begin{align*}
  \nabla \cdot \vv = \prt{v_x}x + \prt{v_y}y + \prt{v_z}z
\end{align*}
It tells us whether our point is a source or a sink: are the fields pointing in and out.\gline{}
Let's consider the problem with the vector function:
\begin{align*}
  \vv &= \frac{\arhat}{r^2}
\end{align*}
We first are going to sketch it, and then compute it's divergence. 
\gimage{image.png}{5cm}
When we try and compute the divergence, we get the following:
\begin{align*}
  \vv &= \frac{\arhat}{r^2}\\
  \vv&= \frac{\pvec{r_i}}{(r_i^2)^{3/2}}\\
  \prt{\vv}{r_i}\cpr{\frac{r_i/r}{\sqrt{r_i^2+\cdots}}^3}&=\frac{1}{r^3}-\frac{3r_i^2}{r^5}
\end{align*}
\begin{align*}
  \sum_{i=1}^s\prt{\vv}{r_i} = \frac{3}{r^5}\left[1-\frac{1}{r^2}(r_i^2+\cdots)\right]=0
\end{align*}
\subsection*{Curl}
Consider the curl, which represents the \textbf{tendency to rotate}:
\begin{align*}
  \nabla\times \vv = \begin{bmatrix}
    \hat x& \hat y&\hat z\\\prt{}x&\prt{}y&\prt{}z\\v_x&v_y&v_z
  \end{bmatrix}
\end{align*}
\newcommand{\del}{\nabla}
Here are some simple product rules for the $\nabla$ operator:
\begin{enumerate}[itemsep=1pt]
  \item $\nabla(fg) = f(\nabla g) + g(\nabla f)$
  \item $\nabla(\vA \cdot \vB)= \vA\times(\nabla\vB)+\vB\times(\nabla\times\vA)+(\vA\cdot\nabla)\vB+(\vB\cdot\nabla)\vA$
  \item $\nabla(f\vA)=f(\nabla\cdot \vA)+\vA\cdot(\nabla f)$
  \item $\del\cdot(\vA\times\vB)=\vB\cdot(\del\times\vA)-\vA\times(\nabla\times\vB)$
  \item $\del\times(f\vA)=f(\nabla\times\vA)-\vA\times(\del f)$
  \item $\del\times(\vA\times\vB)=(\vB\cdot\del)\vA-(\vA\times\del)\vB+\vA(\del\times\vB)-\vB(\del\cdot\vA)$
\end{enumerate}
\subsection*{Second Derivatives}
Now, there are a couple of things we should know for second derivatives:
\begin{enumerate}
  \item Laplacian (divergence of a gradient)
  \begin{align*}
    \del\cdot(\del T)=\prtd{T}{x}+\prtd{T}y + \prtd{T}z
  \end{align*}
  Note that people also write it this way (laplacian of a vector):
  \begin{align*}
    \del^2\vv &= \del^2v_x\hat{x}+\del^2v_y \hat t + \del^2 v_z \hat z\\
    &=\pvec{\del^2 v_x, \del^2v_y, \del^2 v_z}
  \end{align*}
  \item Now we consider the curl of a gradient, which is always zero:
  \begin{align*}
    \del\times(\del T)=0
  \end{align*}
  \item Gradient of a divergence
  \begin{align*}
    \del(\del\times\vv)
  \end{align*}
  \item Divergence of a curl
  \begin{align*}
    \del\times(\del\times\vv)=0
  \end{align*}
  \item Curl of a curl
  \begin{align*}
    \del\times(\del\times\vv)=\del(\del\times\vv)-\del^2\vv
  \end{align*}
\end{enumerate}

\section*{Feb 13th}
\subsection*{Integral types}
\begin{align*}
  \text{Line Integral: }&\int_{\vec a}^{\vec b}\vv \cdot d\vec \ell\\
  &d\vec\ell = dx\hat x + dy \hat y + dx\hat z\\\\
  \text{Surface Integral}&\int_\text{surface}\vv \cdot d\vec a\\
  &\text{for closed surface: outward is positive direction}
\end{align*}
For example, take the surace integral of the bottom of a cube of length 2, along the bottom axis (which runs from (0,0) to (2,2)). OUr field vector is given by $\vv =\pvec{2xz, (x+2), y(x^2-3)}$. \gline{}
\begin{align*}
  &\int_0^2\int_0^2 \pvec{2xz, (x+2), y(x^2-3)}\cdot \pvec{0, 0, -1}\,dxdy\\
  &= \int_0^2\int_0^2 -y(z^2 - 3)\\
  &= \int_0^2 -y2(z^2 - 3) dy\\
  &= -2\frac{1}2y^2(z^2 - 3)\Big|_0^2\\
  &= -12
\end{align*}
\subsection*{Volume Integral}
\begin{align*}
  \int T \,&d\tau\\
  &d\tau=dxdydz
\end{align*}
\subsection*{Fundamental Theorems}
\begin{align*}
  \text{Gradient:} \int_{\vec a}^{\vec b}(\del T)\cdot d\vec\ell = T(\vec b) - T(\vec a)
\end{align*}
\subsection*{Divergence}
\begin{align*}
  &\int_v (\del\cdot \vv)d\tau = \iint_S \vv\cdot d\vec a &d\vec\ell= dx\hat x + dy\hat y +dz\hat z\\
  &\text{closed surface }d\vec a \text{ is out.}
\end{align*}
\subsection*{Curl}
\begin{align*}
  \underline{\int_S(\del\times \vv)\times d\vec a} = \oint_P \vv\cdot d\vec\ell\\
\end{align*}
\subsection*{Stoke's Theorem}
\begin{align*}
  \int_s(\del\times \vv)\cdot d\vec{a} = \oint \vv\cdot d\vec \ell
\end{align*}
\subsection*{Spherical}
Here's how we do spherical coordinates:
\gimage{image2.png}{5cm}
Here, we have that:
\begin{align*}
  r: 0&\to\infty &z = r\cos\theta\\
  \phi: 0&\to 2\pi &y=r\sin\theta\sin\phi\\
  \theta: 0&\to \pi &z=r\sin\theta\cos\phi
\end{align*}
Then, when we write $d\vec \ell$, we have that:
\begin{align*}
  d\vec\ell = d\ell_r \hat r + d \ell_\theta\hat\theta + d\ell_\phi\hat\phi
\end{align*}
Note that $d \vec a$ depends on surface. Constant - gives direction. The other two are changing.\gline{}
For example, a surface in the $xy$ plane, wehave that:
\begin{align*}
  \theta &\text{ constant }=\frac{\pi}2\\
  d\vec a &= d\ell_r d\ell_\phi\hat\theta\\
  &= dr\, r\sin\theta d\phi\hat\theta\\
  &= r\,dr\,d\theta\hat\theta 
\end{align*}
\subsection*{Dirac Delta Function}
Observe the dirac delta function:
\begin{align*}
  \delta(x) &= \begin{cases}
    \infty,&x = 0\\
    0 &x \ne 0
  \end{cases}\\
  \int_{-\infty}^{\infty}\delta(x) \,dx &= 1\\
  \int_{-\infty}^{\infty}f(x)\delta(x) \,dx &= f(0)\\
  \int_{-\infty}^{\infty}f(x)\delta(x - a) \,dx &= f(a)
\end{align*}
\pagebreak
\subsection*{Vector Fields}
\subsubsection*{Theorem 1}
\[\del \times \vec F = 0\quad \vec F = -\del V\]
\subsubsection*{Theorem 2}
\[\del\cdot\vec F = 0\quad\vec F = \del\times\vec A\]

\subsection*{Coulomb's Law}
\begin{align*}
  \vec F &= \frac{1}{4\pi\varepsilon_0}\frac{q_0 Q}{n^2}\hat{\mathbf n}\\
  \vec E &= \frac{\vec F}{Q}\\
  \vec E &=\frac{1}{4\pi \varepsilon_0}\sum_{i=1}^{n}\frac{q_i\hat {\mathbf{n}_i}^2}{n_i^2}
\end{align*}
\subsection*{Continuous Charge Distribution}
\begin{align*}
  \vec E (\vec r) = \frac{1}{4\pi\varepsilon_0}\int\frac{\hat{\mathbf{n}}}{n^2}\,dq
\end{align*}
\begin{adjustwidth}{14em}{14em}
  \begin{align*}
    &\text{line}\quad &dq = \lambda \;d\ell '\\
    &\text{surface}\quad &sq = \sigma \;da'\\
    &\text{volume}\quad &dq = e\;d\tau'
  \end{align*}
\end{adjustwidth}

\end{document}